この章では、本研究で用いるプログラミング言語であるF\texttt{\#}の記法について説明する。
F\texttt{\#}は2005年にマイクロソフト社によって開発されたマルチパラダイム言語である。
関数型プログラムを中心に設計されており、関数の組み合わせによって処理を記述することを特徴とする。
また、F\texttt{\#}は関数型プログラミングに加えて、手続き型プログラミングやオブジェクト指向プログラミングもサポートしている。
マルチパラダイム言語とは、このように複数のプログラミングパラダイムを併せ持つ言語のことを指す。
\section{関数定義}
F\texttt{\#} では、let を用いて関数を定義する。
関数は引数を空白で区切って記述し、関数本体の最後の式を使用して戻り値と型を決定する。
\par
また、F\texttt{\#}ではすべての関数は値とみなす。
これにより、処理内容を引数として外部から指定することが可能となり、柔軟なプログラムの記述が可能である。
この性質により、処理の抽象化や再利用が容易になる。
本研究では、時間を引数として座標や角度を計算する関数を定義し、これらをアニメーションの描画処理に利用している。
これにより、時間に応じて変化する図形の動きを簡潔に記述している。

\section{クラス定義}
F\texttt{\#}では、typeを用いてクラスを定義することができる。
クラスを用いることで、データとそれに対する処理を一つにまとめて定義することが可能である。
クラスには、メンバーとして関数を定義することができ、これらはメソッドと呼ばれる。メソッドはmemberを用いて定義される。
本研究では、図形やアニメーションの描画処理を管理するためにクラスを用いている。
例えば、直線や円弧、テキストなどのアニメーションは、それぞれメソッドとして定義されており、これらはHTMLクラスのメソッドとして実装されている。
